% Method section, written paragraph by paragraph. Notation follows Fig. 2 (make_method.py):
%   image I, query x, phrases \phi(x), candidate masks m_v (v = 1..N), mask token \hat{x}_v, selection score s_v,
%   answer v^\star. Citation keys as in intro.tex.
\section{Method}
\label{sec:method}

\paragraph{Problem formulation.}
The input is an image $I\in\mathbb{R}^{H\times W\times 3}$ and a text query $x$, for example ``locate the handle of
the mug''. The query refers to at most one region, which may be a whole object or a part of one. The output is a
binary mask $m\in\{0,1\}^{H\times W}$ of that region, or the empty answer $\varnothing$ when nothing in the image
matches the query. Most segmentation MLLMs meet this goal by generating $m$, decoding it from a hidden state of the
language model. We cast the task as selection instead (Fig.~\ref{fig:method}). A frozen segmenter proposes a pool
of $N$ candidate masks $\mathcal{M}(I,x)=\{m_v\}_{v=1}^{N}$, the model assigns each candidate a score $s_v(I,x)$, and
the answer is the best-scoring candidate,
\begin{equation}
  v^{\star}=\arg\max_{v\in\{1,\dots,N\}} s_v(I,x), \qquad
  \hat{m}=\begin{cases} m_{v^{\star}} & \text{if } s_{v^{\star}}\geq\tau,\\ \varnothing & \text{otherwise,}\end{cases}
  \label{eq:selection}
\end{equation}
where the threshold $\tau$ is set on validation and is needed only on benchmarks where a query can have no target.
Under this formulation the extent of the answer is fixed by the segmenter, and the model has one thing left to
learn: which candidate the query means. The following subsections describe how the pool $\mathcal{M}$ is built
(Sec.~\ref{sec:pool}), how each candidate becomes a token the language model can read (Sec.~\ref{sec:tokens}), how
the scores $s_v$ are computed (Sec.~\ref{sec:selection}), and how the model is trained (Sec.~\ref{sec:training}).

\subsection{Query-conditioned candidate pool}
\label{sec:pool}

Selection can only succeed if the right mask is in the pool. No single way of prompting a segmenter guarantees
this, so we prompt a frozen SAM\,3~\cite{sam3} in two ways and take the union,
\begin{equation}
  \mathcal{M}(I,x)=\mathcal{M}_{g}(I)\,\cup\,\mathcal{M}_{t}(I,x).
  \label{eq:pool}
\end{equation}

\paragraph{Query-free masks $\mathcal{M}_g$.}
We place a $G\times G$ grid of points on the image and prompt SAM\,3 with each point. This gives masks at all
scales, from whole objects to small parts, and it also covers regions that have no common name. These masks do not
depend on the query. Many of them are near-copies of each other, so we drop a mask when its IoU with a mask we
already kept exceeds $\delta_g$. From the rest we keep the $N_g$ masks with the highest stability score.

\paragraph{Query-conditioned masks $\mathcal{M}_t$.}
A point grid can miss a thin or small part, or split it into several pieces. So we also ask SAM\,3 for the query
itself. We turn the query $x$ into a short list of noun phrases $\phi(x)$. For a part query such as ``the handle of
the mug'', $\phi(x)$ holds the part phrase and the object phrase (``the handle of the mug'', ``the mug''). For an
object query it holds the object phrase, and for any other query it holds the query text. We prompt SAM\,3 with
each phrase and add its $k$ highest-scoring instances to the pool, skipping an instance when its IoU with one
already added exceeds $\delta_t$. For each of these masks we store its rank $r_v\in\{1,\dots,k\}$ among the
instances of its phrase and its SAM\,3 score $\sigma_v$; Sec.~\ref{sec:tokens} gives both to the model.

The pool therefore has
\begin{equation}
  N=|\mathcal{M}(I,x)|\ \le\ N_g+k\,|\phi(x)|
\end{equation}
masks. Both sources matter: on part queries, the best mask in the pool comes from the point grid in $52$--$59\%$ of
the cases and from the text prompts in the rest (Tab.~\ref{tab:coverage}).

\paragraph{Containment tree.}
Masks in the pool overlap. The handle lies inside the mug, and the mug may lie inside a larger region. We record
this structure as a tree over the pool. The parent of a mask $m_v$ is the smallest larger mask that contains at
least a fraction $\rho$ of it,
\begin{equation}
  \mathrm{parent}(v)=\arg\min_{u:\,|m_u|>|m_v|}\,|m_u| \quad \text{s.t.}\quad |m_u\cap m_v|\ \ge\ \rho\,|m_v| .
  \label{eq:tree}
\end{equation}
A mask with no such container is a child of the root, which stands for the whole image. The tree needs no labels
and no training. We use it twice: the position of a mask in the tree is part of its token (Sec.~\ref{sec:tokens}),
and the tree relation between two masks sets a learned bias on the attention between their tokens
(Sec.~\ref{sec:selection}).

\paragraph{Settings.}
We use $G=32$, $N_g=256$, $k=3$, $\delta_g=0.9$, $\delta_t=0.95$ and $\rho=0.8$ in all experiments, which gives
$N\approx260$ candidates per query.
